\Article{9}
Der Vergleich der Gleichungen (7), (9) zeigt, dass die Hilfsgrößen \(A,B,C\) etc. die Werte der Unbestimmten \(x,y,z\) etc. sind, die diesen Werten der Unbestimmten \(\xi,\eta,\zeta\) etc. entsprechen: \(\xi=\mathfrak A,\ \eta=\mathfrak B,\ \zeta=\mathfrak C\) etc. Daraus ist ersichtlich, dass
\[
\left.\begin{aligned}\relax
[aa]A+[ab]B+[ac]C+\text{etc.}&=\mathfrak A\\
[ab]A+[bb]B+[bc]C+\text{etc.}&=\mathfrak B\\
[ac]A+[bc]B+[cc]C+\text{etc.}&=\mathfrak C
\end{aligned}\right\}\ (12)
\]
etc. gilt. Multiplizieren wir daher die Gleichungen (10) jeweils mit \(\frac ap,\ \frac{a'}{p'},\ \frac{a''}{p''}\) etc.\(\cdot\) und addieren sie, so erhalten wir
\[
\left.\begin{aligned}
\mathfrak A&=a\varepsilon+a'\varepsilon'+a''\varepsilon''+\text{etc.}\\
\text{und auf ganz gleiche Weise}\quad&\\
\mathfrak B&=b\varepsilon+b'\varepsilon'+b''\varepsilon''+\text{etc.}\\
\mathfrak C&=c\varepsilon+c'\varepsilon'+c''\varepsilon''+\text{etc.}
\end{aligned}\right\}\ (13)
\]
etc. Da nun \(\mathfrak A\) der Wert der Funktion \(X\) ist, wenn für \(v,v',v''\) etc. die beobachteten Werte eingesetzt werden, sieht man leicht: Wendet man auf jene jeweils die Korrekturen \(-\varepsilon,\ -\varepsilon',\ -\varepsilon''\) etc. an, so erhält die Funktion \(X\) den Wert \(0\); ebenso werden die Funktionen
\SourceLeaf{072}
\(Y,Z\) etc. dadurch auf einen verschwindenden Wert gebracht. Ebenso folgt aus der Gleichung (11), dass \(K\) der Wert der Funktion \(u\) ist, der sich aus derselben Einsetzung ergibt.

Die Anwendung der Korrekturen \(-\varepsilon,-\varepsilon',-\varepsilon''\) etc. auf die Beobachtungen nennen wir \textit{Ausgleichung der Beobachtungen}; damit sind wir offenbar zu einem Schluss von größter Bedeutung geführt worden: Die auf die angegebene Weise ausgeglichenen Beobachtungen genügen sämtlichen Bedingungsgleichungen genau und liefern für jede beliebig von den Beobachtungen abhängige Größe gerade denjenigen Wert, der sich aus der geeignetsten Kombination der unveränderten Beobachtungen ergeben würde. Da es also unmöglich ist, die Fehler selbst \(e,e',e''\) etc. aus den Bedingungsgleichungen zu ermitteln, deren Anzahl dazu nicht ausreicht, haben wir wenigstens die \textit{plausibelsten Fehler} erhalten; so dürfen wir die Größen \(\varepsilon,\varepsilon',\varepsilon''\) etc. bezeichnen.

\Article{10}
Da die Anzahl der Beobachtungen als größer als die Anzahl der Bedingungsgleichungen vorausgesetzt wird, lassen sich neben dem System der plausibelsten Korrekturen \(-\varepsilon,-\varepsilon',-\varepsilon''\) etc. unendlich viele andere finden, die den Bedingungsgleichungen genügen; es lohnt sich zu untersuchen, wie sich diese zu jenem verhalten. Es mögen daher \(-E,-E',-E''\) etc. ein solches, vom plausibelsten verschiedenes System bilden; dann haben wir
\[
\begin{aligned}
aE+a'E'+a''E''+\text{etc.}&=\mathfrak A\\
bE+b'E'+b''E''+\text{etc.}&=\mathfrak B\\
cE+c'E'+c''E''+\text{etc.}&=\mathfrak C
\end{aligned}
\]
etc. Multiplizieren wir diese Gleichungen jeweils mit \(A,B,C\) etc. und addieren sie, so erhalten wir mithilfe der Gleichungen (10)
\[
p\varepsilon E+p'\varepsilon'E'+p''\varepsilon''E''+\text{etc.}=A\mathfrak A+B\mathfrak B+C\mathfrak C+\text{etc.}
\]
Auf ganz gleiche Weise liefern aber die Gleichungen (13)
\[
p\varepsilon\varepsilon+p'\varepsilon'\varepsilon'+p''\varepsilon''\varepsilon''+\text{etc.}=A\mathfrak A+B\mathfrak B+C\mathfrak C+\text{etc.}\quad(14)
\]
Aus der Verbindung dieser beiden Gleichungen ergibt sich leicht
\[
\begin{aligned}
&pEE+p'E'E'+p''E''E''+\text{etc.}\\
&=p\varepsilon\varepsilon+p'\varepsilon'\varepsilon'+p''\varepsilon''\varepsilon''+\text{etc.}
\\&\qquad+p(E-\varepsilon)^2+p'(E'-\varepsilon')^2+p''(E''-\varepsilon'')^2+\text{etc.}
\end{aligned}
\]
\SourceLeaf{073}
Die Summe \(pEE+p'E'E'+p''E''E''+\text{etc.}\) wird also notwendig \textit{größer} sein als die Summe \(p\varepsilon\varepsilon+p'\varepsilon'\varepsilon'+p''\varepsilon''\varepsilon''+\text{etc.}\). Dies lässt sich als folgender Satz aussprechen:

\textsc{Satz.} \textit{Die Summe der jeweils mit den Gewichten der Beobachtungen multiplizierten Quadrate der Korrekturen, durch welche die Beobachtungen mit den Bedingungsgleichungen in Einklang gebracht werden können, wird ein Minimum, wenn die plausibelsten Korrekturen angenommen werden.}

Dies ist gerade das Prinzip der kleinsten Quadrate, aus dem sich auch die Gleichungen (12), (10) leicht unmittelbar ableiten lassen. Für diese kleinste Summe, die wir im Folgenden mit \(S\) bezeichnen wollen, liefert uns übrigens die Gleichung (14) den Ausdruck \(\mathfrak A A+B\mathfrak B+C\mathfrak C+\text{etc.}\)

\clearpage
